% GENERATED FILE. Edit canonical lessons, then run npm run generate:books.
% canonical-source-hash: fnv1a64:f98842ac98bbd1dd
% canonical-lessons: SA-C102-sthali, SA-C102-casakah, SA-C102-camasah, SA-C102-culli, SA-C102-vartika

\chapter{Plate, Cup, Spoon, Stove, Candle}
\label{ch:sa-plate-cup-spoon-stove-candle}
\hlchaptermodality{102}

\begin{chapteropening}
\textbf{By the end of this chapter:} \emph{I can say a plate, a cup, a spoon, a stove and a candle.}

Complete the last lesson of chapter 102: i can say a plate, a cup, a spoon, a stove and a candle.
\end{chapteropening}

\section[\texorpdfstring{\sk{स्थाली} (sthālī)}{स्थाली (sthālī)}]{\sk{स्थाली} (sthālī) — a plate, a pot}

\label{lesson:SA-C102-sthali}



\begin{quote}
Before the new one: say the Sanskrit for glasses, then the Sanskrit for a bag.
\end{quote}

\subsection*{You'll want to know: \sk{स्थाली}}
\textbf{\sk{स्थाली}} — \emph{sthālī} — \textquotedblleft{}a plate, a pot\textquotedblright{}.

\begin{grammarlens}[title={What you've built}]
One more word for this chapter.
\end{grammarlens}

\subsection*{Your turn}
\begin{itemize}
\raggedright
  \item \emph{Say it:} \emph{sthālī}
  \item \emph{Say it:} \emph{sthālī}, once more
  \item \emph{Say it:} say \emph{sthālī}
  \item \emph{Recall:} say the Sanskrit for glasses, then the Sanskrit for a bag, then say \emph{sthālī} again
\end{itemize}

\begin{tcolorbox}[breakable,colback=teal!4,colframe=teal!35!black,title={Before you move on}]
What does \sk{स्थाली} mean? (\textquotedblleft{}A plate, a pot\textquotedblright{}.) Say it once more.
\end{tcolorbox}
\section[\texorpdfstring{\sk{चषकः} (caṣakaḥ)}{चषकः (caṣakaḥ)}]{\sk{चषकः} (caṣakaḥ) — a cup}

\label{lesson:SA-C102-casakah}



\begin{quote}
Before the new one: say the Sanskrit for a bag, then the Sanskrit for a plate.
\end{quote}

\subsection*{You'll want to know: \sk{चषकः}}
\textbf{\sk{चषकः}} — \emph{caṣakaḥ} — \textquotedblleft{}a cup\textquotedblright{}.

\begin{grammarlens}[title={What you've built}]
One more word for this chapter.
\end{grammarlens}

\subsection*{Your turn}
\begin{itemize}
\raggedright
  \item \emph{Say it:} \emph{caṣakaḥ}
  \item \emph{Say it:} \emph{caṣakaḥ}, once more
  \item \emph{Say it:} say \emph{caṣakaḥ}
  \item \emph{Recall:} say the Sanskrit for a bag, then the Sanskrit for a plate, then say \emph{caṣakaḥ} again
\end{itemize}

\begin{tcolorbox}[breakable,colback=teal!4,colframe=teal!35!black,title={Before you move on}]
What does \sk{चषकः} mean? (\textquotedblleft{}A cup\textquotedblright{}.) Say it once more.
\end{tcolorbox}
\section[\texorpdfstring{\sk{चमसः} (camasaḥ)}{चमसः (camasaḥ)}]{\sk{चमसः} (camasaḥ) — a spoon}

\label{lesson:SA-C102-camasah}



\begin{quote}
Before the new one: say the Sanskrit for a plate, then the Sanskrit for a cup.
\end{quote}

\subsection*{You'll want to know: \sk{चमसः}}
\textbf{\sk{चमसः}} — \emph{camasaḥ} — \textquotedblleft{}a spoon\textquotedblright{}.

\begin{grammarlens}[title={What you've built}]
One more word for this chapter.
\end{grammarlens}

\subsection*{Your turn}
\begin{itemize}
\raggedright
  \item \emph{Say it:} \emph{camasaḥ}
  \item \emph{Say it:} \emph{camasaḥ}, once more
  \item \emph{Say it:} say \emph{camasaḥ}
  \item \emph{Recall:} say the Sanskrit for a plate, then the Sanskrit for a cup, then say \emph{camasaḥ} again
\end{itemize}

\begin{tcolorbox}[breakable,colback=teal!4,colframe=teal!35!black,title={Before you move on}]
What does \sk{चमसः} mean? (\textquotedblleft{}A spoon\textquotedblright{}.) Say it once more.
\end{tcolorbox}
\section[\texorpdfstring{\sk{चुल्ली} (cullī)}{चुल्ली (cullī)}]{\sk{चुल्ली} (cullī) — a stove}

\label{lesson:SA-C102-culli}



\begin{quote}
Before the new one: say the Sanskrit for a cup, then the Sanskrit for a spoon.
\end{quote}

\subsection*{You'll want to know: \sk{चुल्ली}}
\textbf{\sk{चुल्ली}} — \emph{cullī} — \textquotedblleft{}a stove\textquotedblright{}.

\begin{grammarlens}[title={What you've built}]
One more word for this chapter.
\end{grammarlens}

\subsection*{Your turn}
\begin{itemize}
\raggedright
  \item \emph{Say it:} \emph{cullī}
  \item \emph{Say it:} \emph{cullī}, once more
  \item \emph{Say it:} say \emph{cullī}
  \item \emph{Recall:} say the Sanskrit for a cup, then the Sanskrit for a spoon, then say \emph{cullī} again
\end{itemize}

\begin{tcolorbox}[breakable,colback=teal!4,colframe=teal!35!black,title={Before you move on}]
What does \sk{चुल्ली} mean? (\textquotedblleft{}A stove\textquotedblright{}.) Say it once more.
\end{tcolorbox}
\section[\texorpdfstring{\sk{वर्तिका} (vartikā)}{वर्तिका (vartikā)}]{\sk{वर्तिका} (vartikā) — a candle, a wick}

\label{lesson:SA-C102-vartika}



\begin{quote}
Before the new one: say the Sanskrit for a spoon, then the Sanskrit for a stove.
\end{quote}

\subsection*{You'll want to know: \sk{वर्तिका}}
\textbf{\sk{वर्तिका}} — \emph{vartikā} — \textquotedblleft{}a candle, a wick\textquotedblright{}.

\begin{grammarlens}[title={What you've built}]
One more word for this chapter.
\end{grammarlens}

\subsection*{Your turn}
\begin{itemize}
\raggedright
  \item \emph{Say it:} \emph{vartikā}
  \item \emph{Say it:} \emph{vartikā}, once more
  \item \emph{Say it:} say \emph{vartikā}
  \item \emph{Recall:} say the Sanskrit for a spoon, then the Sanskrit for a stove, then say \emph{vartikā} again
\end{itemize}

\begin{tcolorbox}[breakable,colback=teal!4,colframe=teal!35!black,title={Before you move on}]
What does \sk{वर्तिका} mean? (\textquotedblleft{}A candle, a wick\textquotedblright{}.) Say it once more.
\end{tcolorbox}
